\subsection{闭算子、可闭算子与稠定算子}

在本节中， K 表示一个交换拓扑域（TG, III, p. 54）。

定义 2. --- 设 E 与 F 是 K 上的拓扑向量空间（EVT, I, p. 1，定义 1）。设 \(u \in \mathcal{P}(\mathrm{E};\mathrm{F})\) 是从 E 到 F 的部分算子。

若 u 的定义域在 E 中稠密，就称 u 是稠定算子。

若 u 的图像在拓扑向量空间 E × F 中是闭的，就称 u 是闭算子。若 u 有一个闭扩张，就称 u 是可闭的。

设 E 、 F 与 G 是 K 上的拓扑向量空间。稠定算子 \(u \in \mathcal{P}(\mathrm{E};\mathrm{F})\) 的每个扩张都是稠定的。此外，若 \(v \in \mathcal{L}(\mathrm{E};\mathrm{F})\) ，则 \(u + v\) 是稠定算子。若 \(v: F \to G\) （相应地， \(w: G \to E\) ）是拓扑向量空间的同构，则 \(v \circ u\) （相应地， \(u \circ w\) ）是稠定算子。

例. --- 设 E 与 F 是 K 上的拓扑向量空间。

\begin{enumerate}
\def\labelenumi{\arabic{enumi})}
\item
  设空间 F 是豪斯多夫的。设 u 是从 E 到 F 中的线性映射。若 u 连续，则部分算子 u 是闭的（TG, I, p. 53，推论 2）。再假设 K 是非离散赋值域，且 E 与 F 是 K 上的可度量化的完备拓扑向量空间。由闭图像定理（EVT, I, p. 19，推论 5），这时由 u 所定义的部分算子是闭的，当且仅当 u 连续。
\item
  设空间 E 是豪斯多夫的。设 v 是从 F 到 E 中的单射连续线性映射。这时部分算子 \(v^{-1} \in \mathcal{P}(\mathrm{E};\mathrm{F})\) （参见 IV, p. 226）是闭的，因为它的图像是 v 的图像在从 \(F \times E\) 到 \(E \times F\) 的、由 \((y,x) \mapsto (x,y)\) 所定义的拓扑向量空间同构下的像，而 v 的图像是闭的。
\end{enumerate}

命题 1. --- 设 E 与 F 是 K 上的拓扑向量空间。从 E 到 F 中的部分算子 u 是可闭的，当且仅当 u 的图像 \(\Gamma_{u}\) 在 E × F 中的闭包是一个泛函图像。这时存在唯一的从 E 到 F 的部分算子 v ，其图像为 \(\overline{\Gamma}_{u}\) ，且它是 u 的最小闭扩张。

若 \(u\) 的图像在 \(\mathrm{E} \times \mathrm{F}\) 中的闭包是泛函图像，则它就是某个部分算子的图像，而后者是 \(u\) 的一个闭扩张，因而它是可闭的。反之，设 \(u \subset w\) ，其中 \(w\) 是闭的。闭包 \(\overline{\Gamma}_u\) （即 \(u\) 的图像在 \(\mathrm{E} \times \mathrm{F}\) 中的闭包）包含在 \(\Gamma_w\) 中，因而 \(\overline{\Gamma}_u\) 是泛函图像。

最后的断言由下述事实得出：若 w 是 u 的一个闭扩张，则 w 的图像包含 \(\overline{\Gamma}_{u}\) 。

定义 3. --- 设 E 与 F 是 K 上的拓扑向量空间。设 u 是从 E 到 F 中的可闭算子。图像为 \(\overline{\Gamma}_{u}\) 的闭算子称为 \(u\) 的闭包，记作 \(\overline{u}\) 。

注. --- 设 E 与 F 是 K 上的拓扑向量空间。设 u 是从 E 到 F 的可闭算子。u 的闭包的定义域包含在 u 的定义域在 E 中的闭包内。一般而言，二者并不相等（IV, p. 344 的习题 1, b)）。

若 \(u \in \mathcal{P}(\mathrm{E};\mathrm{F})\) 是可闭的且 \(\operatorname{dom}(u) = \operatorname{E}\) ，则 \(u = \overline{u}\) 是闭的，因为这时 \(\operatorname{dom}(\overline{u}) = \operatorname{dom}(u)\) 。

命题 2. --- 设 K = R ，并设 E 与 F 是 R 上的拓扑向量空间。设 u 是从 E 到 F 中的部分算子。则 u 是稠定的（相应地，是闭的、是可闭的），当且仅当 \(u_{(\mathbf{C})}\) 是稠定的（相应地，是闭的、是可闭的）。

设 \(u\) 的定义域在 E 中稠密。 \(\mathrm{E}_{(\mathbf{C})}\) 中 0 的每个邻域都包含形如 \(\mathrm{V} + i\mathrm{V}\) 的一个邻域（EVT, II, p. 65），其中 V 是 E 中 0 的一个邻域，因而它包含 \(u_{(\mathbf{C})}\) 的定义域中的一个元素；所以这一部分算子是稠定的。反过来也成立，因为从 \(\mathrm{E}_{(\mathbf{C})}\) 到 E 中的、把 \(x + iy\) 映到 \(x\) （对所有 \((x,y) \in \mathrm{E} \times \mathrm{E}\) ）的映射是连续且满射的。

我们将 \(u_{(\mathbf{C})}\) 的图像等同于 \(u\) 的图像的复化拓扑向量空间。于是，它在 \(\mathrm{E}_{(\mathbf{C})} \times \mathrm{E}_{(\mathbf{C})}\) 中是闭的，如果 \(u\) 的图像在 \(\mathrm{E} \times \mathrm{E}\) 中是闭的。反之，我们有 \(\Gamma_u = \Gamma_{u_{(\mathbf{C})}} \cap (\mathrm{E} \times \mathrm{E})\) （在 \(\mathrm{E}_{(\mathbf{C})} \times \mathrm{E}_{(\mathbf{C})}\) 中）；由于 \(\mathrm{E} \times \mathrm{E}\) 在 \(\mathrm{E}_{(\mathbf{C})} \times \mathrm{E}_{(\mathbf{C})}\) 中是闭的，故部分算子 \(u\) 在 \(u_{(\mathbf{C})}\) 闭时是闭的。

从 E 到 F 中的部分算子 \(v\) 是 \(u\) 的扩张，当且仅当 \(v_{(\mathbf{C})}\) 是 \(u_{(\mathbf{C})}\) 的扩张；因此，部分算子 \(u_{(\mathbf{C})}\) 在 \(u\) 可闭时是可闭的。反之，若 \(u_{(\mathbf{C})}\) 可闭，则部分算子 \(u\) 也可闭，因为 \(\overline{\Gamma_u} = \overline{\Gamma_{u_{(\mathbf{C})}}} \cap (\mathrm{E} \times \mathrm{E})\) 这时是泛函图像（命题 1）。

引理 1. — 设 E 、 F 与 G 是 K 上的拓扑向量空间。设 u 是从 E 到 F 的闭算子。

\begin{enumerate}
\def\labelenumi{\alph{enumi})}
\item
  对每个 \(v \in \mathcal{L}(\mathrm{E};\mathrm{F})\) ，部分算子 \(u + v\) 是闭的。
\item
  对每个 \(v \in \mathcal{L}(\mathrm{G};\mathrm{E})\) ，部分算子 \(u \circ v\) 是闭的。
\end{enumerate}

我们来证明 \(a\) )。设 \(\gamma\) 是映射 \((x,y)\mapsto (x,y - v(x))\) ，它从 \(\mathrm{E}\times \mathrm{F}\) 到其自身，并且是连续的。对每个 \((x,y)\in \mathrm{E}\times \mathrm{F}\) ，我们有 \(\gamma (x,y)\in \Gamma_u\) 当且仅当 \(x\in \operatorname {dom}(u)\) 且 \(y = u(x) + v(x)\) ，亦即 \(\overline{\gamma}^{1}(\Gamma_{u}) = \Gamma_{u + v}\) 。断言得证。

我们来证明 \(b\) )。映射 \(\eta = (v, 1_{\mathrm{F}})\) 从 \(\mathbf{G} \times \mathbf{F}\) 到 \(\mathbf{E} \times \mathbf{F}\) 中是连续的；对每个 \((z, y) \in \mathbf{G} \times \mathbf{F}\) ，我们有 \(\eta(z, y) = (v(z), y)\) ，从而 \(\overline{\eta}(\Gamma_u) = \Gamma_{u \circ v}\) 。断言得证。

设 E 与 F 是 K 上的拓扑向量空间，且 F 是豪斯多夫的。设 \(a \in \mathrm{K}\) 。若 \(u \in \mathcal{P}(\mathrm{E};\mathrm{F})\) 是可闭的，则 \(au\) 也是可闭的。若 \(a \neq 0\) ，我们有 \(\overline{au} = a\overline{u}\) ，且 \(u\) 是闭的当且仅当 \(au\) 是闭的。若 \(a = 0\) ，则 \(au\) 的闭包是 \(0_{\overline{\operatorname{dom}(u)}}\) ，而 \(a\overline{u}\) 等于 \(0_{\operatorname{dom}(\overline{u})}\) ；因此可能出现这样的情形： \(u\) 是闭的，而 \(au\) 不是闭的。

命题 3. --- 设 E 与 F 是 K 上的拓扑向量空间，其中 F 是豪斯多夫的。设 u 是从 E 到 F 的闭部分算子。u 的核是 E 的一个闭向量子空间。

事实上，u 的核是 E × F 的闭子空间 \(\Gamma_{u} \cap (\mathrm{E} \times \{0\})\) 在从 E 到 E × F 中的连续线性映射 \(x \mapsto (x, 0)\) 下的原像。

在本节余下部分，我们假设 K 是非离散赋值域。

定义 4. --- 设 E 与 F 是 K 上的赋范空间，u 是从 E 到 F 的部分算子。对 dom(u) 中的 x ，我们记

\[
\| x \| _ {u} = (\| x \| _ {\mathrm{E}} ^ {2} + \| u (x) \| _ {\mathrm{F}} ^ {2}) ^ {1 / 2}.
\]

如此定义的映射 \(x \mapsto \|x\|_{u}\) 是 \(\text{dom}(u)\) 上的一个范数。我们把这样得到的赋范空间记作 \(E_{u}\) 。

注. --- 设 E 与 F 是 K 上的赋范空间，u 是从 E 到 F 中的部分算子。

\begin{enumerate}
\def\labelenumi{\arabic{enumi})}
\item
  \(E_{u}\) 到 E 中的典则单射是连续的，因为我们有 \(\|x\|\leqslant\|x\|_{u}\) （对所有 \(x\in E\) ）。特别地，\(\operatorname{dom}(u)\) 的在 E 中闭的每个子空间在 \(E_{u}\) 中也是闭的。
\item
  若 E 与 F 是希尔伯特空间，则空间 \(E_{u}\) 是准希尔伯特空间，因为 \(E_{u}\) 上的范数来自下述正埃尔米特形式
\end{enumerate}

\[
(x, y) \mapsto (x \mid y) _ {u} = \langle x \mid y \rangle + \langle u (x) \mid u (y) \rangle
\]

（它定义在 dom(u) 上）。

命题 4. --- 设 E 与 F 是巴拿赫空间（相应地，希尔伯特空间），u 是从 E 到 F 中的部分算子。则 u 是闭的，当且仅当赋范空间 \(E_{u}\) 是巴拿赫空间（相应地，希尔伯特空间）。

只须处理巴拿赫空间的情形。 \(E_{u}\) 的范数，是经由双射线性映射 \((x,y)\mapsto x\) （它把 \(\Gamma_{u}\) 映到 dom(u) 中）作结构迁移，从下述范数得到的：即限制到子空间 \(\Gamma_{u}\) （赋以范数 \((x,y)\mapsto(\|x\|_{\mathrm{E}}^{2}+\|y\|_{\mathrm{F}}^{2})^{1/2}\) ，它由巴拿赫空间 \(E\oplus F\) 上的范数限制而得）上所得到的范数。因此，空间 \(E_{u}\) 是巴拿赫空间，当且仅当子空间 \(\Gamma_{u}\) （在 \(E\oplus F\) 中）是闭的。

若 \(u\) 与 \(v\) 分别是从 E 到 F 与从 F 到 G 的部分算子，且若 \(u\) 是闭的，则部分算子 \(v \circ u\) 一般说来不是闭的，即使 \(v\) 是连续的（IV, p. 344 的习题 1, c)）。尽管如此，我们有如下的充分条件：

引理 2. --- 设 E 、 F 与 G 是 K 上的赋范空间，其中 F 是巴拿赫空间。设 u 是从 E 到 F 的闭部分算子，并设 \(v \in \mathcal{L}(\mathrm{F};\mathrm{G})\) 。如果存在 \(\mathrm{C} \in \mathbf{R}_{+}\) 使得

\[
\| u (x) \| \leqslant C (\| x \| + \| (v \circ u) (x) \|)
\]

对所有 \(x \in \text{dom}(v \circ u) = \text{dom}(u)\) 成立；也就是说，若从 \(x \mapsto u(x)\) 到 F 中的线性映射 \(E_{v \circ u}\) 是连续的，则 \(v \circ u\) 是闭的。

记 \(w = v \circ u\) 。设 \((x_n, w(x_n))_{n \in \mathbf{N}}\) 是 \(w\) 的图像中的一个序列，它在 \(E \times G\) 中收敛。设 \(x\) 是序列 \((x_n)\) 的极限。假设表明，序列 \((u(x_n))_{n \in \mathbf{N}}\) 这时是 \(F\) 中的一个柯西序列；它收敛到元素 \(y\) ，该元素属于 \(F\) 。由于 \(u\) 是闭的，我们有 \(x \in \text{dom}(u)\) 与 \(y = u(x)\) 。于是 \(w(x_n) = v(u(x_n))\) 趋于 \(v(y) = v(u(x))\) ，因为 \(v\) 连续；因此 \(w\) 的图像是闭的。

设 u 是巴拿赫空间 E 上的闭部分算子， F 是 \(\text{dom}(u)\) 的一个子空间。若 F 在巴拿赫空间 \(E_{u}\) 中稠密，则 u 到 F 上的约化是可闭的，且其闭包等于 u 。这时称 F 为 u 的一个核。

